This chapter interprets the results presented in Chapter~5 in
relation to the two research questions defined in Chapter~1,
positions the proposed approach against related work, and reflects
on the limitations and threats to validity of the study. The
following sections address RQ1 and RQ2 in turn, before discussing
the findings in relation to existing work, their practical
implications, the limitations of the study, and directions for
future work.

\section{RQ1: LLM-Based Assessment as a Complement to Rule-Based Assessment}

RQ1 asked to what extent LLM-based assessment can complement
rule-based FAIRness assessment by capturing qualitative documentation
characteristics of research software repositories. The results
presented in Chapter~5 support a qualified but consistent answer:
the LLM-based component adds assessment coverage beyond the structural
indicators evaluated by the rule-based component by addressing
qualitative characteristics of README documentation that deterministic
repository checks do not capture.

The clearest evidence for this complementarity is the difference in
assessment scope. The rule-based component evaluates the presence of
structural FAIR-related evidence such as persistent identifiers,
metadata files, licenses, dependency specifications, README files,
and citation information. The LLM-based component instead evaluates
installation clarity, usage clarity, purpose description, dependency
documentation, examples or tutorials, and contribution guidance.
These dimensions address different properties of the repository and
are therefore retained separately in the hybrid profile.

The descriptive relationship between the structural and documentation
components provides additional evidence for this distinction. For the
primary Llama-based profile, the rule-based and documentation totals
showed only a weak and statistically non-significant association
across the 30 ground truth repositories
($\rho=0.175$, $p=0.356$). This result does not establish statistical
independence between the two components, but it is consistent with
them capturing different aspects of the assessed repositories rather
than providing redundant measurements of the same characteristics.

The example profile presented in Section~4.6 illustrates this
difference concretely. \texttt{astropy/astropy} satisfies all six
rule-based indicators (6/6) while receiving a comparatively low
documentation score from the primary Llama configuration (4/12).
A purely rule-based assessment therefore indicates that the repository
satisfies all six structural indicators evaluated in this work,
whereas the documentation component identifies limited concrete
installation and usage guidance within the README evidence supplied
to the model. This illustrates the motivation for the hybrid design:
the presence of structural FAIR-related artefacts does not by itself
describe the clarity and completeness of documentation available to
a prospective software user.

The baseline comparison in Chapter~5 provides further evidence of
the difference in scope. The six documentation-quality dimensions
evaluated by the LLM-based component are not represented in howfairis,
which focuses on structural properties such as licensing, citation,
software-registry information, and checklist-style evidence. The
F-UJI pilot points in the same direction from a different perspective.
When applied directly to five GitHub repositories with substantially
different structural and documentation profiles, F-UJI returned the
same overall FAIRness score for all five. Within this limited pilot,
it therefore did not distinguish between repositories whose
software-oriented profiles differed substantially, consistent with
its primary design for assessing research data objects.

Complementarity should, however, be distinguished from reliability.
The two evaluated LLM configurations produced substantially different
levels of agreement with the human annotations. Against the primary
annotator, Llama achieved a macro-average weighted Kappa of 0.329,
whereas Qwen achieved 0.612; against the second annotator, the
corresponding values were 0.275 and 0.557. Human--human agreement was
0.709. The documentation component should therefore not be interpreted
as a uniformly reliable substitute for human judgement merely because
it extends the assessment scope. Instead, the results show that
qualitative documentation assessment can provide complementary
information while its reliability remains dependent on the model
configuration and the documentation dimension being assessed.

RQ1 can therefore be answered in terms of assessment coverage rather
than replacement. The LLM-based component complements deterministic
FAIR-related checks by evaluating qualitative documentation
characteristics that are otherwise absent from the implemented
structural assessment. The hybrid profile makes these additional
characteristics visible without collapsing them into the structural
scores or a single aggregate FAIRness value. The reliability with
which the qualitative dimensions can be assessed is a separate
question and is addressed by RQ2.

\section{RQ2: Reliability of LLM-Based Assessment and Its Failure Modes}

RQ2 asked how reliably LLM-based methods can assess qualitative documentation
dimensions of research software repositories, and what failure modes affect
their assessments. The two-model evaluation shows that there is
no single reliability level that characterises LLM-based assessment
in this setting. Agreement depends substantially on the model
configuration and the documentation dimension being assessed, while
the disagreement patterns additionally reveal limitations associated
with the evidence available to the model.

The clearest model-level difference is observed in the aggregate
agreement results. Against the primary annotator, Llama achieved a
macro-average weighted Kappa of 0.329, whereas Qwen achieved 0.612.
Against the second annotator, the corresponding values were 0.275
and 0.557. Human--human agreement across the six documentation
dimensions was 0.709. The human value provides useful context for
interpreting model--human agreement, but it should not be treated as
a formal upper bound on model performance. The difference between
the two LLM configurations nevertheless demonstrates that conclusions
about the reliability of qualitative documentation assessment cannot
be generalised from the behaviour of a single model.

The dimension-level results reinforce this conclusion. Against the
primary annotator, Qwen achieved its strongest agreement for
installation (D1, $\kappa=0.799$), purpose
(D3, $\kappa=0.789$), and contribution information
(D6, $\kappa=0.759$). Usage documentation reached
$\kappa=0.516$, while dependencies (D4, $\kappa=0.487$) and examples
(D5, $\kappa=0.320$) remained more difficult. Llama showed a
different pattern, with substantially weaker aggregate agreement and
particularly low agreement for usage, purpose, and examples.

Purpose description (D3) provides the clearest illustration of model
sensitivity. Llama achieved a negative weighted Kappa
($\kappa=-0.071$) against the primary annotator despite 86.7\% exact
agreement. The primary human scores were strongly concentrated at
score 2, with 28 of 30 repositories receiving that score, making
Kappa sensitive to the small number of disagreements. Qwen, however,
agreed with the primary annotator on 29 of 30 repositories and
achieved $\kappa=0.789$. Against the second annotator, Qwen achieved
100\% exact agreement and $\kappa=1.000$. The weak Llama result for
D3 was therefore not reproduced by the second model configuration.
This indicates that the Llama D3 result should not be interpreted as
evidence that purpose description is inherently unsuitable for
LLM-based assessment.

At the same time, the stronger Qwen results did not remove
dimension-specific failure patterns. Examples and tutorials (D5)
remained difficult under both configurations. Relative to the primary
annotator, both Llama and Qwen assigned lower D5 scores in 17 of the
30 repositories. Qwen never assigned a higher D5 score than the
primary annotator and achieved only $\kappa=0.320$ for this dimension.
Against the second annotator, the pattern was even more pronounced:
Qwen assigned lower D5 scores in 26 of 30 repositories and achieved
$\kappa=0.133$.

Dependency documentation (D4) was an additional source of systematic
disagreement for Qwen. The model assigned a lower D4 score than the
primary annotator in 13 of 30 repositories and never assigned a
higher score, despite its substantially stronger aggregate agreement
relative to Llama. These results show that improvement at the
model level does not imply uniformly reliable assessment across all
documentation dimensions.

One plausible contributor to some of the underscoring is the
README-only evidence policy. The automated assessment was restricted
to the supplied README text and did not follow external documentation
links. This is particularly relevant for dimensions such as usage
(D2) and examples (D5), for which repositories may direct users to
external guides, tutorials, or example galleries. Human annotators
could account for references to such resources according to the
annotation procedure, whereas the automated assessment was restricted
to the README text supplied to the model. The observed disagreement
is therefore partly shaped by the evidence boundary imposed by the
assessment design. However, the two-model results also show that
evidence scope alone cannot explain all disagreement, since the
models produced substantially different agreement levels while
operating under the same README-only policy.

The pilot-exclusion sensitivity analysis provides additional evidence
that the model-level findings were not primarily driven by the
repositories used during pilot-phase refinement. Llama's
macro-average weighted Kappa increased from 0.329 on the complete
30-repository corpus to 0.373 after excluding the 10 pilot
repositories. Qwen remained almost unchanged, from 0.612 to 0.615.
The stability of the Qwen result in particular indicates that its
observed agreement was not primarily attributable to the original
pilot overlap.

Model-reported confidence did not provide a dependable reliability
signal under either configuration. For Llama, confidence was not
monotonically associated with observed correctness: MEDIUM-confidence
predictions achieved higher exact agreement and lower MAE than
HIGH-confidence predictions overall. Qwen exhibited a different
limitation by assigning HIGH confidence to all 180 dimension-level
assessments. Confidence calibration across levels could therefore not
be evaluated for Qwen, and the uniform confidence label did not
distinguish dimensions with substantially different observed
agreement. For example, Qwen exact agreement ranged from 43.3\% for
D5 to 96.7\% for D3 despite every assessment being labelled HIGH
confidence. Self-reported confidence should therefore not be treated
as a calibrated probability of correctness in the present framework.

README length likewise did not explain the overall disagreement
patterns. At the repository level, the relationship between effective
README length and model--human score difference was weak and
non-significant for both Llama ($\rho=-0.169$, $p=0.371$) and Qwen
($\rho=-0.143$, $p=0.451$). An exploratory dimension-level analysis
for Qwen identified a negative association for D4
($\rho=-0.392$, $p=0.032$), but this result was obtained from a
secondary analysis across six dimensions on only 30 repositories and
without multiple-comparison correction. It therefore does not provide
sufficient evidence for a general README-length effect.

Taken together, the answer to RQ2 is that LLM-based qualitative
documentation assessment can achieve substantial agreement with human
judgement for some model--dimension combinations, but its reliability
is model-dependent, dimension-dependent, and sensitive to the
available evidence. The contrast between Llama and Qwen demonstrates
that poor agreement observed for one model should not automatically
be attributed to the assessment task itself, while the recurring
difficulties for D4 and D5 show that stronger aggregate model
performance does not eliminate criterion-specific weaknesses.
Self-reported confidence also does not provide a validated mechanism
for identifying unreliable assessments. LLM-generated documentation
scores should therefore be interpreted together with criterion-level
evidence and empirical reliability results rather than treated as
uniformly dependable automated judgements.

\section{Comparison with Related Work}

The results of this thesis can be situated against two strands of
related work introduced in Chapter~2: automated LLM-based FAIR
assessment, represented most directly by Fair-Way, and the broader
LLM-as-judge literature on using language models as automated
evaluators.

Fair-Way \cite{sharma2025fairway} is the closest related system,
applying LLMs to automate FAIR assessment against the FAIRsFAIR DOAM
metrics for research datasets. Its authors report that GPT-4o, the
strongest model evaluated, reaches an exact match accuracy of 0.905
against their benchmark, while smaller and open-source models show
lower performance. These results are not directly comparable to the
agreement values reported in this thesis because the studies evaluate
different constructs, use different scoring scales and metrics, and
target different resource types. Fair-Way evaluates FAIR-related
dataset metadata, whereas the LLM-based component developed here
assesses qualitative README documentation for research software using
an ordinal rubric.

Nevertheless, the two-model evaluation in this thesis is consistent
with Fair-Way's broader observation that model choice can materially
affect LLM-based assessment performance. Under the same six
documentation dimensions and README-only evidence policy, Llama
achieved a macro-average weighted Kappa of 0.329 against the primary
annotator, whereas Qwen achieved 0.612. The contrast was especially
pronounced for purpose description (D3), where the corresponding
values were $-0.071$ and 0.789. These results provide direct evidence
within the present task that reliability depends partly on the model
configuration rather than solely on the assessment rubric or evidence
source.

The present evaluation remains considerably narrower than a systematic
model benchmark. Only two model configurations were evaluated, and
they differ in model family, scale, execution environment, and output
handling. The results therefore do not establish a general performance
ordering between LLM families or show that one model class is
intrinsically more suitable for FAIRness assessment. Instead, they
demonstrate that conclusions obtained from a single model
configuration may not generalise even when the assessment dimensions,
evidence policy, and scoring rubric are held constant.

More substantively, this thesis differs from Fair-Way in scope and
architecture rather than only in observed model performance. Fair-Way
targets research datasets and evaluates structured and semi-structured
metadata using an LLM-based pipeline. Its authors identify hybrid
approaches combining LLM assessment with rule-based verification as
a direction for further work, motivated in part by the limitations
of relying exclusively on generative model outputs. The framework
developed in this thesis investigates such a hybrid design in the
context of research software repositories: deterministic rules are
used for structural FAIR-related indicators, while LLMs are restricted
to qualitative characteristics of unstructured README documentation.

The complementarity finding discussed in Section~6.1 provides
empirical support for this separation of responsibilities. The
rule-based component evaluates repository properties for which
explicit deterministic evidence can be defined, whereas the LLM-based
component addresses documentation characteristics requiring semantic
interpretation. The weak and non-significant relationship observed
between the structural total and the primary Llama-based documentation
total is consistent with the two components contributing different
information to the resulting profile rather than acting as duplicate
measures of the same repository properties.

The results also relate to the broader LLM-as-judge literature.
Zheng et al.\ \cite{zheng2023llmjudge} report that strong language
models can achieve high agreement with human evaluators in some
open-ended evaluation settings while also identifying systematic
sources of unreliability such as position bias, verbosity bias, and
self-enhancement bias. The present study differs substantially in
task design: the models evaluate repository documentation against a
fixed six-dimension ordinal rubric rather than comparing free-form
responses. Nevertheless, both studies illustrate that aggregate
agreement alone is insufficient for characterising the reliability
of an LLM-based evaluator.

In the present experiment, no statistically significant
repository-level verbosity effect was observed for either evaluated
model configuration. The relationship between effective README length
and model--human score difference was weak for both Llama
($\rho=-0.169$, $p=0.371$) and Qwen
($\rho=-0.143$, $p=0.451$). An exploratory Qwen result for dependency
documentation (D4) showed a negative dimension-level association, but
the small sample and multiple dimension-level comparisons do not
support treating this as evidence of a general verbosity effect.

Instead, the dominant reliability findings in this thesis were
model- and dimension-specific. Examples and tutorials (D5) remained
difficult under both model configurations, while Qwen also showed
systematic underscoring for dependency documentation (D4). In
addition, the README-only evidence policy limits what the models can
assess when repositories place relevant material in external
documentation. These findings represent task-specific limitations of
the assessment setup rather than evidence for a general bias of
LLM-based evaluators.

The confidence results provide a further connection to the broader
problem of trusting automated LLM judgements. Neither evaluated
configuration produced a self-reported confidence signal that could
be treated as a calibrated estimate of correctness. Llama confidence
was not monotonically associated with observed agreement, while Qwen
assigned HIGH confidence to every dimension-level assessment. This
reinforces the importance of validating reliability empirically rather
than inferring it from the model's own expression of certainty.

Overall, the comparison with related work positions the contribution
of this thesis primarily in its hybrid architecture and
software-specific assessment scope rather than in direct numerical
performance against existing systems. Fair-Way demonstrates the
potential of LLM-based FAIR assessment and the importance of model
choice in a research-data setting, while the present work evaluates
a combination of deterministic structural assessment and semantic
documentation assessment for research software. The two-model results
further show that model sensitivity must be considered explicitly
when interpreting the reliability of the qualitative component.

\section{Implications}

The findings of this thesis carry several practical implications for
how a hybrid FAIRness assessment framework of this kind should be
used and interpreted, beyond the specific numerical results reported
in Chapter~5.

The most direct implication concerns how the two assessment
components should be consumed together. Because the rule-based and
LLM-based components were designed to capture different types of
evidence, and the observed scores showed only a weak relationship
(Section~6.1), a hybrid profile is most useful when read as two
separate signals rather than as a single quality indicator. A
repository with a high rule-based score should not be assumed to
have well-documented, easily reusable software, and a repository
with a low rule-based score is not necessarily poorly documented.
The example presented in Section~4.6 makes this concrete: a
repository can satisfy every structural indicator evaluated in
this work while still leaving a new user with limited concrete
guidance on how to install and use the software. Presenting
the two dimensions separately, rather than combining them into an
overall FAIRness score, allows this kind of gap to remain visible
rather than being averaged away.

A second implication concerns how the LLM-based scores themselves
should be interpreted. The two-model evaluation shows that their
reliability cannot be represented by a single general performance
level. Llama achieved a macro-average weighted Kappa of 0.329 against
the primary annotator, whereas Qwen achieved 0.612 under the same
documentation dimensions and evidence policy. At the same time,
agreement varied considerably across individual dimensions for both
models. LLM-generated documentation scores should therefore be
interpreted with reference to the model configuration and the
dimension being assessed rather than treated as uniformly reliable
automated judgements.

This distinction has practical consequences for human review. The
stronger Qwen agreement demonstrates that qualitative documentation
assessment can be substantially more reliable under one model
configuration than another, but persistent dimension-specific
disagreements remain. In particular, examples and tutorials (D5)
showed systematic underscoring under both configurations, while
dependency documentation (D4) showed pronounced underscoring for
Qwen. Low scores in such dimensions should therefore be interpreted
together with the available criterion-level evidence, especially when
relevant documentation may be located outside the README supplied to
the model.

A related implication concerns the confidence labels produced
alongside each score. Neither evaluated configuration provided a
self-reported confidence signal that could be used as a validated
estimate of assessment reliability. For Llama, HIGH-confidence
predictions did not consistently outperform MEDIUM-confidence
predictions. Qwen exhibited a different limitation by assigning HIGH
confidence to all 180 dimension-level assessments, despite substantial
variation in observed agreement across dimensions. The confidence
field should therefore be treated as contextual model output rather
than as a calibrated probability of correctness or an automatic basis
for deciding which assessments require human review.

A more appropriate review strategy would use the assessment evidence
and empirically observed dimension-level reliability alongside the
predicted score. Dimensions with recurring disagreement patterns can
be treated more cautiously, while the deterministic structural
indicators remain directly traceable to explicit repository evidence.
This preserves the practical benefit of automated semantic assessment
without hiding the uncertainty identified through the evaluation.

Finally, the baseline comparison suggests a practical positioning
for the proposed approach relative to existing tools rather than a
claim of outright superiority. Where howfairis provides fast,
purely structural checks and F-UJI is designed for research data
objects rather than software repositories, the LLM-based component
of this framework adds a distinct assessment
layer --- documentation quality as experienced by a prospective user
--- that is not directly represented by either tool in the
comparisons conducted in this thesis. This suggests the two kinds of
tools are better viewed as complementary than as competitors: an
organisation or repository maintainer seeking fast, deterministic
structural FAIR-related checks may use a tool such as howfairis,
while the qualitative documentation assessment introduced
in this thesis addresses a different, additional question that such
tools do not answer.

\section{Limitations and Threats to Validity}

This section discusses the limitations of the study and the threats
to validity that should be considered when interpreting its results.

\textbf{Limited model scope.}
The LLM-based component was evaluated using two model configurations:
a locally hosted Llama~3.1 8B model and
\texttt{Qwen3.6-35B-A3B-MLX-8bit}, accessed through a
university-hosted API. The substantial difference in agreement between
the two configurations demonstrates that model choice affects the
observed reliability of the assessment. However, evaluating two
configurations is still insufficient to establish model-independent
properties of LLM-based documentation assessment or a general
performance ordering across model families. The models also differ in
architecture, scale, execution environment, and output handling.
The findings should therefore be interpreted as characterising the
two evaluated configurations under the assessment setup used in this
thesis rather than as estimates of LLM performance in general.

\textbf{Qwen output-recovery procedure.}
The initial Qwen evaluation successfully produced valid assessments
for 23 of the 30 repositories, while seven responses failed output
parsing because the returned content was incomplete or malformed.
For the recovery run, the maximum output-token allowance was increased
from 4{,}096 to 16{,}000 tokens and the output instruction was amended
to request that each reasoning field contain only one short sentence.
The 23 successful original assessments were retained, while only the
seven failed repositories were evaluated again under the modified
output constraint. The substantive D1--D6 rubric, README evidence,
input-length limit, score scale, and temperature remained unchanged.
Nevertheless, the final Qwen result set was produced under two
slightly different output-generation configurations. This introduces
a minor protocol variation that should be considered when interpreting
the Qwen agreement results, even though the modification concerned
response generation and formatting rather than the assessment
criteria themselves.

\textbf{Input truncation strategy.} README content exceeding
12{,}000 characters was truncated using a fixed sequential cutoff
rather than a token-based or section-aware strategy, as described in
Chapter~4. Within the 30-repository ground truth corpus, only one
README exceeded this threshold, and manual inspection of the
truncated portion showed no loss of relevant D1--D6 evidence in that
particular case. This single instance does not establish that fixed
character-based truncation is safe in general; a section-aware
strategy that preserves installation, usage, dependency, example, and
contribution content specifically would be more robust for longer
or differently structured READMEs than the corpus evaluated here.

\textbf{Representativeness of the ground truth corpus.} The mean
rule-based score across the 30-repository ground truth corpus (4.30
out of 6) was somewhat higher than the mean across the full
188-repository dataset (3.86 out of 6). Because the ground truth
corpus was constructed to maximise diversity across research domains
and programming-language ecosystems rather than sampled at random
from the full dataset, it may not be fully representative of the
broader population of repositories to which the rule-based scorer
was subsequently applied. Agreement and validation results reported
against the ground truth corpus should accordingly be interpreted as
characterising performance on a deliberately diverse but non-random
sample, rather than as an unbiased estimate of performance across
the full dataset.

\textbf{Partial exclusion logging.} During dataset construction, the
keyword-based exclusion step was logged explicitly, recording the
specific repositories removed for containing predefined
non-software-related keywords. Repositories excluded for being forks,
archived, or inactive since before the collection cutoff date were
filtered during the initial collection stage but were not separately
logged. Consequently, a complete, itemised record of every excluded
repository and its specific exclusion reason is not available for
the full 198-to-188 filtering process, limiting the transparency of
this step relative to the ideal of a fully auditable exclusion trail.

\textbf{Ground truth correction process.} Corrections made to the
ground truth annotations during validation --- including the
dependency-file schema correction and the license-annotation
correction for \texttt{rust-ml/linfa} discussed in Chapter~5 --- were
applied directly to the annotation records rather than being tracked
in a separate change log. While each correction was traceable to a
specific, documented cause and was applied consistently using the
finalised annotation guidelines, the absence of a separate audit
trail means that the sequence of corrections cannot be independently
reconstructed from the released annotation files alone.

\textbf{Scale of the ground truth corpus.} The 30 manually annotated
repositories enable the agreement and validation analyses reported in
this thesis, but Kappa and correlation estimates computed on a sample
of this size carry non-trivial uncertainty, as illustrated by the
prevalence-effect result observed for D3 in Section~6.2. A larger
annotated corpus would allow for more precise estimates and a more
fine-grained analysis of agreement across research domains and
programming languages.

\textbf{README-only evidence scope.} The LLM-based assessment was
restricted to the README content supplied to the model and did not
retrieve or inspect documentation available through external links.
This restriction ensured that every repository was evaluated against
a consistent and reproducible evidence source, but it also limits the
construct being measured. In particular, repositories that place
usage guides, tutorials, example galleries, or other documentation
outside the README may receive lower scores even when such resources
are available to users.
The qualitative error analysis in Section~6.2 indicates that this
evidence restriction is particularly relevant to dimensions such as
D2 and D5, for which documentation is frequently referenced through
external guides, tutorials, or example resources. However, the
substantial differences observed between the two model configurations
under the same README-only policy show that evidence scope alone does
not explain the observed disagreement. The reported LLM scores should
therefore be interpreted as assessments of documentation available
directly from the README rather than of the complete documentation
ecosystem of a repository.

\textbf{Subjectivity of qualitative ground truth.} The documentation
dimensions D1--D6 require human judgement about the completeness and
clarity of repository documentation. Although the substantial
agreement between the two human annotators
($\kappa=0.709$) provides evidence that the rubric can be applied
with reasonable consistency, disagreement remained for individual
dimensions and repositories. The manual annotations should therefore
not be treated as an error-free absolute ground truth, but as a
reference assessment produced according to the operationalised rubric
defined in Chapter~3. This uncertainty should be considered when
interpreting LLM--human disagreement, since some differences may
reflect ambiguity in the assessment task itself rather than model
error alone.

\textbf{Scope of the baseline comparison.} The comparison with
existing FAIR assessment tools was limited in scope. While howfairis
was applied to the ground truth repositories, the F-UJI analysis was
conducted as a five-repository pilot because the tool is primarily
designed for research data objects rather than GitHub software
repositories. The identical F-UJI scores observed in this pilot are
therefore illustrative of the behaviour encountered for the selected
repositories and should not be interpreted as a comprehensive
benchmark of F-UJI against the proposed framework.

\section{Future Work}

The findings and limitations discussed in this chapter point to
several directions for extending the proposed framework.

\textbf{Extending the multi-model evaluation.}
The evaluation of Llama~3.1 8B and
\texttt{Qwen3.6-35B-A3B-MLX-8bit} demonstrated that model
configuration can substantially affect agreement with human
annotations. Future work should extend this comparison to a broader
set of model families, sizes, and deployment settings while keeping
the documentation rubric and evaluation procedure fixed. A larger
model comparison would help determine which of the observed
dimension-specific patterns recur across models and which are
associated primarily with individual model configurations. It would
also provide a stronger basis for investigating whether particular
documentation dimensions require different prompting strategies or
assessment mechanisms.

\textbf{Extension to other hosting platforms.}
The current implementation is limited to repositories hosted on
GitHub. The framework could be extended to other research-software
hosting and archival platforms, such as Zenodo or Codeberg, which
would require adapting the repository collection and
evidence-extraction components to platform-specific metadata and APIs.

\textbf{Section-aware and token-based input handling.} The current
implementation applies a fixed character-based truncation cutoff to
README content, as described in Chapter~4. Future work could
implement a section-aware preprocessing strategy that identifies and
preserves content relevant to each documentation dimension ---
installation, usage, dependencies, examples, and contribution
guidance --- regardless of its position in the README, together with
a token-based length limit that more directly reflects the model's
actual context capacity.

\textbf{Addressing the external-link limitation.}
The README-only evidence policy prevents the LLM-based component from
evaluating documentation available through external links, which is
particularly relevant to dimensions such as usage guidance and
examples. Future work could investigate retrieving and incorporating
linked documentation --- such as dedicated usage guides, tutorials,
or example galleries --- into the assessment evidence. Such an
extension would make it possible to test directly how much of the
observed disagreement is attributable to the restricted evidence
scope rather than to the model configuration itself. Any retrieval
strategy would need to preserve reproducibility, clearly record the
evidence supplied to the model, and control the additional assessment
cost.

\textbf{Repeated-run stability analysis.} The current evaluation
reports scores from a single inference run per repository. Running
the LLM-based assessment repeatedly for the same repositories would
allow score consistency and dimension-level variation across runs to
be measured directly, complementing the model-reported confidence
analysis in Chapter~5.

\textbf{Scaling and diversifying the ground truth corpus.} Expanding
the manually annotated corpus beyond 30 repositories, with a
sampling strategy closer to the composition of the full dataset,
would reduce the estimation uncertainty discussed in Section~6.5 and
allow agreement and validation results to be reported with greater
confidence, including analysis broken down by research domain and
programming-language ecosystem.

\textbf{Pipeline integration and extended user-facing presentation.}
The current implementation consists of independently executable
components, as described in Chapter~4. Future work could integrate
these components behind a single command-line interface, reducing
the need to invoke each processing stage separately, and package the
framework for straightforward execution against a GitHub repository.
The web-based prototype introduced in Chapter~4 could additionally be
extended with further usability and presentation improvements ---
for example, batch live-assessment of multiple repositories or
export of hybrid profiles in machine-readable form --- building on
the dataset overview, repository-level, and live assessment views
already implemented.

\textbf{Improved uncertainty and review mechanisms.}
Neither evaluated model configuration produced a self-reported
confidence signal that reliably distinguished correct from less
reliable assessments. Future work should therefore investigate
uncertainty measures derived from observable model behaviour rather
than relying solely on self-reported confidence. Possible approaches
include consistency across repeated inference runs, agreement between
multiple model configurations, or evidence-based checks that compare
the generated score with the supporting documentation excerpts.
Such signals could support a human-review mechanism in which
uncertain, unstable, or conflicting assessments are flagged for
inspection while deterministic structural indicators remain fully
automated.